% ICLR 2027 paper template for optimal-cot-depth project
% Theory paper: complexity-theoretic analysis of optimal CoT depth

\documentclass{article}

\usepackage[final]{iclr2027}  % Change to [preprint] for arXiv version

% Essential packages for theory papers
\usepackage{amsmath,amssymb,amsthm}
\usepackage{algorithm}
\usepackage{algorithmic}
\usepackage{graphicx}
\usepackage{booktabs}
\usepackage{hyperref}
\usepackage{cleveref}

% Theorem environments (shared counter)
\newtheorem{theorem}{Theorem}
\newtheorem{lemma}[theorem]{Lemma}
\newtheorem{corollary}[theorem]{Corollary}
\newtheorem{proposition}[theorem]{Proposition}
\theoremstyle{definition}
\newtheorem{definition}[theorem]{Definition}
\newtheorem{assumption}[theorem]{Assumption}
\newtheorem{example}[theorem]{Example}
\theoremstyle{remark}
\newtheorem{remark}[theorem]{Remark}

% Complexity notation
\newcommand{\TC}{\mathrm{TC}^0}
\newcommand{\NC}{\mathrm{NC}^1}
\newcommand{\NP}{\mathrm{NP}}
\newcommand{\coNP}{\mathrm{coNP}}
\newcommand{\PSPACE}{\mathrm{PSPACE}}
\newcommand{\poly}{\mathrm{poly}}

% CoT-specific notation
\newcommand{\CoT}{\mathrm{CoT}}
\newcommand{\depth}{d^*}       % optimal depth function
\newcommand{\noise}{\eta}      % per-step noise rate
\newcommand{\acc}{\mathrm{acc}} % accuracy function

% Standard math
\newcommand{\E}{\mathbb{E}}
\newcommand{\R}{\mathbb{R}}
\newcommand{\N}{\mathbb{N}}
\newcommand{\bigO}{\mathcal{O}}

\title{On the Optimal Depth of Chain-of-Thought:\\A Complexity-Theoretic Analysis}

\author{Anonymous Authors \\
  \texttt{\{anonymous\}@institution.edu}
}

\begin{document}

\maketitle

\begin{abstract}
Recent work reaches contradictory conclusions about Chain-of-Thought (CoT) reasoning depth: some studies show longer chains improve performance, while others demonstrate degradation. We resolve this contradiction through computational complexity theory. We prove that the optimal number of reasoning steps $\depth(\mathcal{C}, n)$ is determined by the complexity class $\mathcal{C}$ of the underlying task and the input size $n$. For tasks in P, optimal depth is $\bigO(1)$---additional steps introduce noise without benefit. For NP tasks, depth scales polynomially with input size. Beyond a complexity-determined threshold, additional steps degrade performance via hallucination cascades at a rate governed by the per-step error probability $\noise$. We formalize this as a noise accumulation model and prove tight bounds for standard complexity classes. Empirical validation across tasks of known complexity confirms our predictions: the observed optimal depth matches the theoretical curve within [TBD]\% across [TBD] models. Our framework explains why each contradictory finding in the literature is correct---for the specific complexity profile of its evaluation tasks.
\end{abstract}

\section{Introduction}
\label{sec:intro}

% TODO: Develop introduction
% Key points:
% 1. CoT has become the dominant reasoning paradigm for LLMs
% 2. Three papers reach contradictory conclusions about CoT length
% 3. We resolve this through complexity theory
% 4. Preview of main results

\section{Related Work}
\label{sec:related}

% TODO: Develop related work
% Subsections:
% - Chain-of-Thought reasoning and scaling
% - Computational complexity of LLM reasoning (connect to reasoning-gaps)
% - Circuit complexity and transformer expressiveness
% - Error accumulation in sequential computation

\section{Formal Framework}
\label{sec:framework}

% TODO: Develop formal framework
% Key definitions:
% - CoT depth vs. CoT length
% - Optimal depth function d*(C, n)
% - Noise accumulation model
% - "Useful" vs. "noise" reasoning steps

\section{Main Results}
\label{sec:results}

% TODO: Develop main results
% Theorems:
% 1. Depth bound for P tasks (O(1))
% 2. Depth scaling for NP tasks (polynomial in n)
% 3. Noise ceiling theorem (depth beyond threshold degrades performance)
% 4. Unification theorem (resolves the contradiction)

\subsection{Optimal Depth for Polynomial-Time Tasks}
\label{sec:p-bound}

\subsection{Optimal Depth for NP Tasks}
\label{sec:np-bound}

\subsection{The Noise Ceiling}
\label{sec:noise-ceiling}

\subsection{Unifying Contradictory Findings}
\label{sec:unification}

\section{Empirical Validation}
\label{sec:experiments}

% TODO: Develop experiments
% - Benchmark suite with tasks of known complexity class
% - Varying CoT depth (1, 2, 4, 8, 16, 32 steps)
% - Multiple models (GPT-4, Claude, open-source)
% - Compare observed optimal depth to predicted d*(C, n)

\section{Discussion}
\label{sec:discussion}

% TODO: Develop discussion
% - Implications for CoT prompt design
% - Connection to reasoning-gaps and verification-complexity
% - Limitations and future work
% - Broader impact: principled CoT depth selection

\bibliography{references}
\bibliographystyle{iclr2027}

\appendix

\section{Full Proofs}
\label{app:proofs}

% TODO: Full proofs of all theorems

\section{Experimental Details}
\label{app:experiments}

% TODO: Full experimental setup, hyperparameters, task specifications

\section{Notation Table}
\label{app:notation}

\begin{table}[h]
\centering
\begin{tabular}{ll}
\toprule
\textbf{Symbol} & \textbf{Meaning} \\
\midrule
$\depth(\mathcal{C}, n)$ & Optimal CoT depth for complexity class $\mathcal{C}$, input size $n$ \\
$\noise$ & Per-step error (noise) probability \\
$\acc(k, \mathcal{C}, n)$ & Accuracy as a function of $k$ reasoning steps \\
$\TC$ & Threshold circuit complexity class \\
$\NC$ & Nick's class (log-depth circuits) \\
\bottomrule
\end{tabular}
\end{table}

\end{document}
